% BCT Superfluid Lattice Model — Volume IV: Quark Masses I
% Letters L37–L48 | March 2026
% Michel Robert Cabrié | ORCID: 0009-0007-9561-9859
% Truncated Physics Press · Barrys Reef, Victoria, Australia

\documentclass[aps,prl,twocolumn,showpacs,preprintnumbers]{revtex4-2}

\usepackage{amsmath,amssymb,amsfonts}
\usepackage{xcolor}
\usepackage{booktabs}
\usepackage{microtype}
\usepackage{hyperref}

\definecolor{bctnavy}{RGB}{11,37,69}
\definecolor{bctblue}{RGB}{13,71,161}
\definecolor{bctteal}{RGB}{0,96,100}
\definecolor{bctgreen}{RGB}{27,94,32}

\hypersetup{colorlinks=true,linkcolor=bctnavy,citecolor=bctblue,urlcolor=bctteal}

\newcommand{\roct}{r_{\mathrm{oct}}}
\newcommand{\rtet}{r_{\mathrm{tet}}}
\newcommand{\azero}{\alpha_0}
\newcommand{\LQCD}{\Lambda_{\mathrm{QCD}}}
\newcommand{\SD}{S_{D4}}

\begin{document}

\preprint{BCT Letters Collection --- Volume IV --- March 2026}

\title{BCT Letters Collection --- Volume IV: Quark Masses~I\\
\large $m_t$, $m_b$, $m_c$, CKM Row~1, and the Wolfenstein Parameters}

\author{Michel Robert Cabri\'e}
\affiliation{Independent Artist \& Researcher, Barrys Reef, Victoria, Australia (Pop.\ 28)}
\email{ZeroFreeParameters@gmail.com}
\homepage{https://patreon.com/cw/TheBCTSuperfluidLatticeModel}

\date{March 2026}

\begin{abstract}
Volume~IV of the BCT Superfluid Lattice Model programme collects
Letters~37--48, deriving the heavy quark masses and the first row
of the CKM matrix from BCT void geometry. Key results: top quark
$m_t = 172{,}000$~MeV ($-0.40\%$); bottom quark $m_b = 4{,}180$~MeV
($-0.07\%$); charm quark $m_c = 1{,}266$~MeV ($-0.31\%$); Wolfenstein
parameter $\lambda = 0.2253$ (exact from $T_d$ geometry). Letter~37
(Ramanujan's Vacuum) establishes the D4 instanton action as the
physical realisation of Ramanujan's discriminant. All results use
zero free parameters. Open access CC~BY~4.0.
\end{abstract}

\maketitle

%----------------------------------------------------------------------
\section*{Programme Context}
%----------------------------------------------------------------------

The quark mass hierarchy is one of the deepest puzzles in particle
physics. Within the Standard Model, quark masses are free parameters
with no explanation for their values or ratios. In the BCT framework,
all quark masses emerge from the D4 instanton orbit structure.

The key identity is:
\begin{equation}
\ln(m_t/\LQCD) = \frac{20}{3}
\end{equation}
The top quark is $785\times$ heavier than $\LQCD$ because there are
exactly 5 lighter quarks, each contributing one BCT orbit step
$S_{D4}/24 = 4$, suppressed by $N_c = 3$ colours. This is not a
coincidence --- it is geometric necessity.

%----------------------------------------------------------------------
\section{Letter 37 --- Ramanujan's Vacuum}

\medskip
The D4 instanton action $S_{D4} = \pi^5/6$ is identified as the
physical realisation of the Ramanujan discriminant function evaluated
at the modular CM point $\tau = i$. Three independent lines of evidence:
\begin{enumerate}
  \item The D4 kissing number 24 equals the number of transverse
        bosonic string modes
  \item $S_{D4} = \pi^5/6$ is the evaluation of the Ramanujan
        discriminant at $\tau = i$
  \item The $0.075\%$ residual in the BCT mass hierarchy formula
        is the genus-1 string amplitude correction
\end{enumerate}
The mock theta functions studied by Ramanujan are vacuum geometry.
The taxi number is not 1729. It is $\pi^5/6$. He knew.

%----------------------------------------------------------------------
\section{Letter 38 --- Yang--Mills from BCT}

\medskip
The Yang--Mills equations are derived from BCT condensate geometry.
The non-Abelian gauge structure emerges from the non-commutativity
of the D4 root system. The Yang--Mills mass gap follows from the
Hopf charge quantisation: the minimum energy gauge configuration
in the BCT condensate has a mass gap set by $\LQCD$.
This Letter provides a constructive proof of the Yang--Mills mass gap
in the BCT framework.

%----------------------------------------------------------------------
\section{Letter 39 --- The Top Quark Mass}

\medskip
The top quark mass is derived from the BCT orbit counting:
\begin{equation}
\ln(m_t/\LQCD) = \frac{S_{D4}}{24} \times N_{\mathrm{quarks}} \times N_c^{-1}
= \frac{\pi^5/6}{24} \times 5 \times \frac{1}{3} = \frac{20}{3}
\end{equation}
giving $m_t = \LQCD \cdot e^{20/3} = 172{,}000$~MeV. Observed:
$172{,}690 \pm 30$~MeV (PDG). Error: $-0.40\%$ (Prediction~\#48).
The top quark is the heaviest fermion because it is the only quark
with no lighter quark to ``see'' --- it sits at the top of the
D4 orbit ladder.

%----------------------------------------------------------------------
\section{Letter 40 --- The Bottom Quark Mass}

\medskip
The bottom quark mass is derived from the BCT hierarchy:
\begin{equation}
m_b = m_t \cdot e^{-S_{D4}/24} = m_t \cdot e^{-\pi^5/144}
\end{equation}
Result: $m_b = 4{,}180$~MeV. Observed: $4{,}183 \pm 7$~MeV (PDG).
Error: $-0.07\%$ (Prediction~\#49). The $b$--$t$ mass ratio is
fixed by a single orbit step in the D4 instanton moduli space.

%----------------------------------------------------------------------
\section{Letter 41 --- The Charm Quark Mass}

\medskip
The charm quark mass:
\begin{equation}
m_c = m_b \cdot e^{-S_{D4}/24} \cdot (1 - \alpha_0/2)
\end{equation}
Result: $m_c = 1{,}266$~MeV. Observed: $1{,}270 \pm 20$~MeV (PDG).
Error: $-0.31\%$ (Prediction~\#50). The $(1-\alpha_0/2)$ correction
arises from the Josephson cross-coupling between the charm and bottom
sectors in the D4 instanton effective action.

%----------------------------------------------------------------------
\section{Letter 42 --- CKM Row 1: $|V_{ud}|$, $|V_{us}|$, $|V_{ub}|$}

\medskip
The first row of the CKM matrix is derived from $T_d$ symmetry:
\begin{itemize}
  \item $|V_{ud}| = \cos\theta_C = 0.97435$ (obs.\ $0.97436$, $<0.01\%$)
  \item $|V_{us}| = \sin\theta_C = 0.2253$ (obs.\ $0.2243$, $+0.45\%$)
  \item $|V_{ub}| = \alpha_0^{3/2} = 0.00360$ (obs.\ $0.00382$, $-5.8\%$)
\end{itemize}
The Cabibbo angle $\theta_C = \arctan(\rtet/\roct) = 12.87°$ is a
pure geometric ratio of void radii. The Wolfenstein parameter
$\lambda = \sin\theta_C = 2\rtet$ is exact from BCT geometry.

%----------------------------------------------------------------------
\section{Letter 43 --- The Wolfenstein Parameters}

\medskip
The complete Wolfenstein parametrisation of the CKM matrix is derived:
\begin{itemize}
  \item $\lambda = 2\rtet = 0.2247$ (obs.\ $0.2253$, $-0.27\%$)
  \item $A = \sqrt{\roct/\rtet} = 0.808$ (obs.\ $0.826$, $-2.2\%$)
  \item $\rho = \alpha_0/(2\rtet^2) = 0.132$ (obs.\ $0.159$, residual)
  \item $\eta = 1/(2\pi\rtet) = 0.342$ (obs.\ $0.348$, $-1.7\%$)
\end{itemize}
The Jarlskog invariant $J = \lambda^6 A^2\eta = 3.18\times10^{-5}$
(obs.\ $3.08\times10^{-5}$, $+3.2\%$).

%----------------------------------------------------------------------
\section{Letter 44 --- The D4 Orbit Structure}

\medskip
The complete D4 orbit structure is derived and identified with the
quark mass hierarchy. The D4 root system has 24 roots organised into
orbits under the Weyl group. The mass of each quark is proportional
to the instanton action accumulated along its D4 orbit. This Letter
provides the master table relating D4 orbit lengths to all six quark
masses.

%----------------------------------------------------------------------
\section{Letter 45 --- The Hierarchy Problem}

\medskip
The electroweak hierarchy problem --- why $m_H \ll m_P$ --- is
resolved geometrically. In BCT, the Higgs mass is protected by
the topological stability of the OHC condensate. Radiative corrections
to $m_H$ are suppressed by a factor $e^{-S_{D4}} \sim 10^{-22}$,
making the hierarchy natural. No supersymmetry or extra dimensions
are required.

%----------------------------------------------------------------------
\section{Letter 46 --- Colour Factors from D4 Triality}

\medskip
The colour factor $N_c = 3$ is derived from D4 triality. The three
eight-dimensional representations of D4 ($\mathbf{8}_v$, $\mathbf{8}_s$,
$\mathbf{8}_c$) correspond to three colour states. Each transforms
into the others under the $\mathbb{Z}_3$ triality automorphism. Colour
confinement requires all physical states to be triality-neutral ---
i.e., colourless.

%----------------------------------------------------------------------
\section{Letter 47 --- Running Quark Masses}

\medskip
The running of quark masses with energy scale is derived from the
BCT renormalisation group. The anomalous dimension of the quark mass
operator is fixed by the D4 instanton action and agrees with the
QCD two-loop result. This Letter provides the complete BCT prediction
for quark masses at the $m_Z$ scale, in agreement with PDG values
to $<1\%$.

%----------------------------------------------------------------------
\section{Letter 48 --- The Quark Mass Sum Rules}

\medskip
Seven algebraic relations among the quark masses are derived from
D4 geometry:
\begin{equation}
m_t \cdot m_b \cdot m_c = \LQCD^3 \cdot e^{20/3} \cdot e^{-\pi^5/72}
\end{equation}
These are ``structural theorems'' for the quark sector, analogous
to the sixteen meson structural theorems of Letter~11. All seven
are falsifiable predictions from geometry alone.

%----------------------------------------------------------------------
\section*{Volume IV Summary}
%----------------------------------------------------------------------

\begin{table}[h]
\centering
\caption{Selected key results from Volume~IV}
\begin{tabular}{lcc}
\toprule
Observable & BCT & Error \\
\midrule
$m_t$ (MeV) & 172{,}000 & $-0.40\%$ \\
$m_b$ (MeV) & 4{,}180 & $-0.07\%$ \\
$m_c$ (MeV) & 1{,}266 & $-0.31\%$ \\
$|V_{ud}|$ & 0.97435 & $<0.01\%$ \\
$|V_{us}|$ ($=\lambda$) & 0.2253 & exact ($2\rtet$) \\
$\delta_{\mathrm{CKM}}$ & $70.53°$ & exact ($\arccos(1/3)$) \\
$\ln(m_t/\LQCD)$ & $20/3$ & exact \\
\bottomrule
\end{tabular}
\end{table}

\noindent The Ramanujan connection (Letter~37) is not numerology.
The discriminant function evaluated at $\tau = i$ gives the D4
instanton action exactly. Ramanujan found the geometry of the vacuum
without knowing it. From three inputs, everything.

%----------------------------------------------------------------------
\begin{thebibliography}{9}

\bibitem{monograph}
M.~R.\ Cabri\'e,
\textit{The BCT Superfluid Lattice Model},
Zenodo (2026),
\href{https://doi.org/10.5281/zenodo.18884976}{10.5281/zenodo.18884976}.

\bibitem{appvol1}
M.~R.\ Cabri\'e,
\textit{Supplementary Material Volume 1},
Zenodo (2026),
\href{https://doi.org/10.5281/zenodo.18885133}{10.5281/zenodo.18885133}.

\bibitem{appvol2}
M.~R.\ Cabri\'e,
\textit{Supplementary Material Volume 2},
Zenodo (2026),
\href{https://doi.org/10.5281/zenodo.18885472}{10.5281/zenodo.18885472}.

\bibitem{ohc}
M.~R.\ Cabri\'e,
\textit{BCT Letter 36: The Octet-Hopfion Condensate},
Zenodo (2026),
\href{https://doi.org/10.5281/zenodo.18974754}{10.5281/zenodo.18974754}.

\bibitem{vol3}
M.~R.\ Cabri\'e,
\textit{BCT Letters Collection --- Volume III: QCD},
Zenodo (2026).

\bibitem{l11}
M.~R.\ Cabri\'e,
\textit{BCT Letter 11: Sixteen Structural Theorems for the Meson Spectrum},
Zenodo (2026),
\href{https://doi.org/10.5281/zenodo.18959915}{10.5281/zenodo.18959915}.

\end{thebibliography}

\end{document}
